\section{Rをつかった回帰分析}
\label{s11_Regression_onR}

\subsection{授業内容}
\subsubsection{科目の中でのこのコマの位置づけ}
\begin{tabular}[t]{ccccccc}
    記述統計[3] & 確率[1] & モデル[5] & 推測・推定[0] & 検定[1] & 実践[3]
\end{tabular}
\subsubsection{概要}
回帰分析や重回帰分析について，基本的な概念や用語・推定値の意味は前回までに習得済みであるが，それが統計環境では実際どのようにして算出されるか，算出された数字はどのような意味があるかについて具体的に理解することを目的とする。座学の中では記号としてしか意味を持たなかったものが，実際の数字として得られるとより理解が深まることにもなる。統計環境はRを用いるが，SPSSそのほかの統計パッケージを用いても同様の出力が得られる。統計環境が変わる可能性もあるので，出力画面の位置で内容を覚えるのではなく，意味的な理解をしておく必要がある。
\subsubsection{コマ主題細目}
\begin{description}

    \item[散布図の描画]線形回帰はその名の通り線形関係をモデルで表現するものであり，非線形な関係や外れ値の存在などがあれば適切な推定値が得られたことにはならない。分析をする前にまずデータの様子を確認するために，散布図の描画は必ず行わなければならない。第\ref{s04_usingR}講の復習を兼ねて，ggplotによる散布図の描画を行う。
    \begin{flushright}
        $\to$ \citet{Kinosady2021}のP.126--138.
    \end{flushright}

    \item[線形回帰の実行]データが読み込まれれば，$lm$関数によって最小二乗法による回帰分析を行うことができる。Rによる関数関係の記述(formula)の方法を理解し，出力結果から適切な読み取りができるよう，どのような用語でどのような出力がなされているかをしっかり確認する。
    \begin{flushright}
        $\to$ \citet{kosugi2018}のP.54--60.
    \end{flushright}

    \item[残差のプロット]$lm$関数から返される結果オブジェクトの中には，予測値$\hat{Y}$がfitted.values，残差がresidualsという変数名で保存されている。これらを取り出したり図示したりすることで，回帰分析の諸特徴を理解する一助となる。数理的な導出は\citet{kosugi2018}のP.127--135にあるが，具体的な数字としてこれを確認し，理解する。
    \begin{flushright}
        $\to$ \citet{kosugi2018}のP.54--60.
    \end{flushright}

    \item[重回帰分析の実行]統計環境において重回帰分析を行うのは，説明項を1つ追加するだけで良いので作業としては容易い。ただし解釈に注意が必要であることを再確認し，また標準化偏回帰係数の算出方法についても理解を進める。また変数を増やす前と増やした後で，重相関係数が増加することを確認する。
    \begin{flushright}
        $\to$ \citet{kosugi2018}のP.63--67.
    \end{flushright}

\end{description}
\subsubsection{キーワード}
\begin{itemize}
    \item ggplot2
    \item lm関数
    \item formula表記
    \item fitted.valueとresiduals
    \item 標準化偏回帰係数
\end{itemize}
\subsection{授業情報}
\paragraph{コマの展開方法}
演習；4号館4FのPCルーム1/2に別れて，実際のPCを使う実習を伴った授業を行う
\paragraph{標準シラバスにおける位置づけ}
科目番号5 心理学統計法；2統計に関する基礎的な知識；C/D/E エクセル，R，SPSS入門
\subsubsection{予習・復習課題}
\paragraph{予習}
第\ref{s04_usingR}講と同じ教室・形式で行われる。RやRStudio，プロジェクトによる管理など基本的なことを改めて教示しないので，前回の復習をかねて空いている時間にPCルームで基本的なR/RStudioの挙動について理解しておくと良い。またRの内部でどのように数字やデータが扱われるかについて，\citet{kosugi2019}のP.24--36などを参考にみておくと良い。
\paragraph{復習}
残差をデータの中で確認し，部分相関係数や偏相関係数，偏回帰係数がどのように算出されるか自分自身で確かめてみると良い。
また，授業時間内に課される課題をRmd形式で提出することが求められる。時間内に学んだ技術を用いて，レポートを提出すること。なおこのレポートの提出は前期の単位認定にあたっての必要条件であり，基準に満たないレポートは再提出を求める。前期試験の前にパスすることが，前期末試験を受験するための前提条件であることを承知しておくこと。

